\documentclass[11pt]{amsart}

\usepackage[letterpaper,margin=1in]{geometry}
\usepackage{amsmath,amssymb,amsthm,mathtools}
\usepackage{microtype}
\usepackage[colorlinks=true,linkcolor=black,citecolor=black,urlcolor=black]{hyperref}

\numberwithin{equation}{section}

\theoremstyle{plain}
\newtheorem{theorem}{Theorem}
\newtheorem{lemma}[theorem]{Lemma}
\newtheorem{proposition}[theorem]{Proposition}
\newtheorem{corollary}[theorem]{Corollary}

\theoremstyle{remark}
\newtheorem{remark}[theorem]{Remark}

\newcommand{\ZZ}{\mathbb Z}
\newcommand{\RR}{\mathbb R}
\newcommand{\CC}{\mathbb C}
\newcommand{\Zn}[1]{\mathbb Z_{#1}}

\newcommand{\ch}[1]{\chi_{#1}}
\newcommand{\wh}[1]{\widehat{#1}}
\newcommand{\psin}{\psi_n}
\newcommand{\psiN}[1]{\psi_{#1}}
\newcommand{\Apsi}{A_{\psin}}
\newcommand{\ApsiN}[1]{A_{\psi_{#1}}}
\newcommand{\Pt}{P_t}
\newcommand{\elln}{\ell_n}
\newcommand{\Lcyc}{L_n}
\newcommand{\coef}{c}

% Notation for the variance and surplus quantities used in the cubic estimates.
% Change the two symbols below to alter S^2 and Q^2 notation globally.
\newcommand{\Ssymbol}{S}
\newcommand{\Qsymbol}{Q}
\newcommand{\Ssq}{\Ssymbol^2}
\newcommand{\Qsq}{\Qsymbol^2}
\newcommand{\Sfour}{\Ssymbol^4}

\newcommand{\TT}{\mathbb T}
\newcommand{\muT}{\pi}
\newcommand{\AT}{A_{\TT}}
\newcommand{\AvgT}[1]{\int_{\TT} #1\,d\muT}
\newcommand{\EntT}[1]{\operatorname{Ent}_{\muT}\!\left(#1\right)}

% Uniform averaging functional. Change \AvgSymbol to alter notation globally.
\newcommand{\AvgSymbol}{\pi}
\newcommand{\Avg}[1]{\AvgSymbol\!\left(#1\right)}

% Entropy notation on the finite cycle. Change \EntMeasure to alter the subscript globally.
\newcommand{\EntMeasure}{\pi}
\newcommand{\Ent}[1]{\operatorname{Ent}_{\EntMeasure}\!\left(#1\right)}
\newcommand{\Energy}{\mathcal E}
\newcommand{\En}{\Energy_n}
\newcommand{\norm}[1]{\left\|#1\right\|}
\newcommand{\abs}[1]{\left|#1\right|}

\newcommand{\Span}{\operatorname{span}}
\newcommand{\eps}{\varepsilon}
\newcommand{\ii}{\mathrm i}
\newcommand{\LSI}{\operatorname{LSI}}
\newcommand{\HC}{\operatorname{HC}}

\title{Sharp log-Sobolev inequalities on finite cyclic groups}

\author[X. Xie]{Xinyuan Xie}
\address{(X.X.) Department of Mathematics, University of California, Irvine, CA 92697, USA}
\email{xinyuax7@uci.edu}



\author[H. Zhang]{Haonan Zhang}
\address{(H.Z.) Department of Mathematics, University of South Carolina}
\email{haonanzhangmath@gmail.com}

\date{\today}

\begin{document}

\begin{abstract}
Let \(\mathbb Z_n\) be the cyclic group equipped with the uniform probability
measure \(\pi\), and let \(A_{\psi_n}\) be the Laplacian
with word length
\[
  \psi_n(k) = \min(k,n-k).
\]
We prove the sharp log-Sobolev inequality
\[
  \textnormal{Ent}_{\pi}(f^2)
  \le 2\pi(f A_{\psi_n} f),
  \qquad f:\mathbb Z_n \to [0,\infty),
\]
for every \(n \ge 4\). The proof is inspired by the recent work of Frank and
Ivanisvili~\cite{FrankIvanisvili2026} on a sharp log-Sobolev inequality for
nearest-neighbor simple random walk. We use their cubic-majorant
reduction, which turns the problem into a 3rd moment estimate; the new point is a blockwise 3rd moment estimate adapted to the word-length multiplier. The same 3rd moment argument also recovers the log-Sobolev inequality for Poisson-semigroup on the circle, first proved by Weissler~\cite{Weissler1980}.  The same sharp inequalities were also
obtained recently by Yao~\cite{Yao2026} by a different method.
\end{abstract}

\maketitle

\keywords{Key words:
Log-Sobolev inequality, hypercontractivity, finite cyclic groups, Poisson semigroups, cubic Sobolev inequalities.}

\section{Introduction}

There are two natural translation-invariant dynamics on the discrete circle, or cyclic group,
\(\Zn{n}\).  The first is the nearest-neighbor simple random walk: from a
point, one jumps to either neighbor with equal probability.  The associated
positive Laplacian is
\[
  \Lcyc (f)(k)=f(k)-\frac{1}{2}(f(k+1)+f(k-1)).
\]
Using the Fourier expansion
\[
f=\sum_{k\in \Zn{n}} \wh f(k)\ch{k}, \qquad  \ch{k}(x):=e^{2\pi \ii kx/n},
\]
we may write $L_n$ as a Fourier multiplier
\[
L_n(f)=\sum_{k\in \Zn{n}}r_n(k)\wh f(k)\ch{k},\qquad r_n(k):=1-\cos(2\pi k/n).
\]

The second can be described by the positive word-length Laplacian $\Apsi$, which is the Fourier
multiplier
\[
  \Apsi f=\sum_{k\in \Zn{n}}\psin(k)\wh f(k)\ch{k},
  \qquad
  \psin(k):=\min(k,n-k).
\]
Equivalently, $-\Apsi$ is the generator of the Poisson semigroup.
Here, one can think about $\psin$ as the word length function with respect to the symmetric generator set $\{-1,1\}$.

In this paper, we shall consider sharp log-Sobolev inequalities, or equivalently hypercontractivity, in these two cases. For classical log-Sobolev inequalities and hypercontractivity, we refer to
\cite{Bonami1970,Gross1975,Beckner1975,Gross2006} and references therein for background.

We shall use the convention that \(\pi\) denotes the uniform probability measure
on \(\Zn{n}\).  Thus, for all $f$ on $\Zn{n}$,
\[
\Avg{f}
=
\frac1n\sum_{x\in\mathbb Z_n}f(x).
\]
For all non-negative $f$ on $\Zn{n}$, we define
\[
\Ent{f}
:=
\Avg{f\log f}
-
\Avg{f}
\log \Avg{f}.
\]


For the nearest-neighbor cycle, the log-Sobolev inequality states that for all $f:\Zn{n}\to \RR$
\begin{equation}\label{lsi1}
    \Ent{f^2}\le c\Avg{f L_n f}.
\end{equation}
The problem of determining the sharp
log-Sobolev constant $c$ goes back to the work of Diaconis and Saloff-Coste on
finite Markov chains \cite{DiaconisSaloffCoste1996}.
After partial progress by Chen and Sheu \cite{ChenSheu2003}, Chen, Liu and
Saloff-Coste \cite{ChenLiuSaloffCoste2008}, and Faust and Fawzi
\cite{FaustFawzi2024}, Frank and Ivanisvili settled the problem for all
\(n\ge4\), proving
that the optimal constant is \(c=2/r_n(1)\)
\cite{FrankIvanisvili2026}. Their key new step is an optimal cubic Sobolev
inequality, which inspired this work.

For the word-length Laplacian, the log-Sobolev inequality reads: for all $f:\Zn{n}\to \RR$
\begin{equation}\label{lsi2}
    \Ent{f^2}\le c\Avg{f \Apsi f}.
\end{equation}
The study of the optimal constant $c$ appeared in works of Beckner, Janson and Jerison \cite{BJJ1983}, Andersson \cite{Andersson2002}, Wolff \cite{Wolff2007},
Junge, Palazuelos, Parcet and Perrin \cite{JPPP2017}.
In two preprints, Yao first \cite{Yao2025} settled the special cases \(n=2^r\) and
\(n=3\cdot2^r\), and then, more recently, all \(n\ge4\) \cite{Yao2026}. 

We remark that the log-Sobolev inequality for the word-length \eqref{lsi2} is stronger than that of the nearest-neighbor simple random walk \eqref{lsi1}.
Indeed, both generators share the same eigenvectors, the characters $\chi_k$'s.
Their eigenvalues, after normalization so that the spectral gap is one, have the following relation:
\begin{equation}\label{eq:eigenvalue comparison}
    \ell_n(k):=\frac{r_n(k)}{r_n(1)}\ge \psi_n(k),\qquad 1\le k\le n-1.
\end{equation}
To see this, write $m=\min\{k,n-k\}\le n/2$, we have
\begin{equation}
    \frac{r_n(k)}{r_n(1)}=\frac{1-\cos(2\pi k/n)}{1-\cos(2\pi/n)}
=\frac{1-\cos(m 2\pi/n)}{1-\cos (2\pi/n)}
\ge m=
\psi_n(k),\qquad 1\le k\le n-1,
\end{equation}
 where the inequality is by induction and the identity
 \[
1-\cos(a+b)-(1-\cos a)-(1-\cos b)
=
4\cos\!\left(\frac{a+b}{2}\right)
\sin\!\left(\frac{a}{2}\right)
\sin\!\left(\frac{b}{2}\right)
\ge 0,
\qquad a,b\ge 0,\ a+b\le \pi.
\]


The main result of this paper is the proof of the sharp log-Sobolev inequality for the word-length case for all $n\ge 4$. The proof is inspired by Frank--Ivanisvili and is different from Yao's method. 

\begin{theorem}\label{thm:main}
For every \(n\ge4\) and every nonnegative \(f:\Zn{n}\to\RR\),
\[
  \Ent{f^2}\le 2\Avg{f\Apsi f}.
\]
The constant \(2\) is sharp.
\end{theorem}

For complex-valued functions, the corresponding form is
\[
  \Ent{|f|^2}
  \le 2\Avg{\overline f\,\Apsi f},
\]
and it follows from the real nonnegative statement by the standard contraction
of Dirichlet forms.  Following standard arguments, one obtains the following
hypercontractivity result of the Poisson semigroup \(\Pt=e^{-t\Apsi}\).  Sufficiency follows from Gross's theorem
\cite{Gross1975} applied to Theorem \ref{thm:main}; necessity follows from the
standard perturbation.

\begin{corollary}[Hypercontractivity]\label{cor:hypercontractivity}
Let \(\Pt=e^{-t\Apsi}\).  For \(1<p\le q<\infty\),
\[
  \norm{\Pt f}_{q}\le \norm{f}_{p}\qquad \text{for all }f:\Zn{n}\to\CC
\]
holds if and only if
\[
  t\ge \frac12\log\frac{q-1}{p-1}.
\]
\end{corollary}

% We also record the standard tensorization consequence.  If
% \(G=\prod_{r=1}^d \Zn{n_r}\), with each \(n_r\ge4\), is equipped with the
% product uniform measure \(\pi_G\), and if
% \[
%   A_G=\sum_{r=1}^d A_{\psi_{n_r}}^{(r)}
% \]
% is the sum of the one-coordinate word-length generators, then entropy
% subadditivity and tensorization of Dirichlet forms give, with the same sharp
% constant,
% \[
%   \operatorname{Ent}_{\pi_G}(f^2)
%   \le 2\pi_G(fA_G f),
%   \qquad f:G\to[0,\infty).
% \]
% We will not use this standard fact below.

\begin{remark}
    It is a standard argument that the above log-Sobolev inequality and hypercontractivity extend to products of $\mathbb{Z}_n$ with the same constants. 
\end{remark}



Our proof of Theorem \ref{thm:main} uses the cubic-majorant reduction idea of Frank--Ivanisvili who observed that 
\[
  2t^2\log t
  \le
  \frac23(t-1)^2(t+2)+(t^2-1),\qquad t\ge 0.
\]
See Lemma \ref{lem:cubic-majorant} below. This reduces the problem to a cubic, or 3rd moment, Sobolev inequality. 
However, the final high-frequency estimate in their nearest-neighbor proof does not adapt directly to the word-length multiplier. In fact, the main obstacle to a literal adaptation of
\cite{FrankIvanisvili2026} can be seen from the proof of their \cite[Lemma 3]{FrankIvanisvili2026}, which uses 
\[
\sum_{j=2}^{n-2}\frac{1}{\ell_n(j)-1}<\frac{3}{2}.
\]

% \[
%   2(\ell_n(j)-1)\simeq 2(j^2-1)
% \]
% in the low-frequency range. Hence the relevant reciprocal sum is comparable
% to \(\sum_{j\ge2}(j^2-1)^{-1}<\infty\).
However, for the word-length multiplier, $\psin(j)-1=j-1$ when $j<n/2$, so the analogous estimate would contain
\[
  \sum_{j=2}^{m}\frac1{j-1}\simeq\log m,
\]
which produces a logarithmic loss.

Our replacement for this final step is a blockwise 3rd moment estimate.  On the Fourier side, we organize
the proof by conjugate pairs of frequencies \(\{j,-j\}\), except in the even case
where the middle frequency \(n/2\) is self-conjugate.  If \(u:\Zn{n}\to \RR\) has mean zero and
\begin{equation}\label{SQ}
     \Ssq=\Avg{u^2},\qquad \Qsq=\Avg{u\Apsi u}-\Ssq,
     \qquad \Ssymbol,\Qsymbol\ge0,
\end{equation}
then we prove
\[
  \abs{\Avg{u^3}}\le 3\Ssq\Qsymbol.
\]
This 3rd moment estimate counts only the cubic Fourier resonances that survive averaging.
For odd cyclic groups these resonances are
\[
  i+j=k,
  \qquad
  i+j+k=n.
\]
For even cyclic groups the same resonances occur, and one must also handle the
self-conjugate middle frequency \(n/2\).  Compared with Frank--Ivanisvili, the main addition is this blockwise 3rd moment estimate: it closes the cubic inequality for the linear word-length spectrum without summing the harmonic reciprocal-surplus weights.


\bigskip

We close the introduction with the continuous-circle counterpart.  Let
\(\TT=\RR/\ZZ\), let \(\muT\) be normalized Haar measure, and let \(\AT\) be
the Poisson generator on \(\TT\),
\[
  \AT (f)=\sum_{k\in\ZZ}\abs{k}\,\wh f(k)\chi_k,
  \qquad
  \chi_k(x)=e^{2\pi\ii kx},
\]
so that
\[
  \AvgT{f\AT f}=\sum_{k\in\ZZ}\abs{k}\,\abs{\wh f(k)}^2.
\]

Weissler \cite{Weissler1980} first proved the sharp Poisson-semigroup log-Sobolev inequality on the circle. 

\begin{theorem}\label{thm:circle}
For every nonnegative trigonometric polynomial \(f\) on \(\TT\),
\[
  \EntT{f^2}\le 2\AvgT{f\AT f}.
\]
% Consequently, by closure, the same inequality holds for nonnegative functions
% in the form domain of \(\AT\) whenever the entropy is finite.  The constant
% \(2\) is sharp.
\end{theorem}

This log-Sobolev inequality for Poisson semigroup can be recovered from the finite word-length
inequalities by a Riemann-sum limit; see
\cite{Yao2026}.  Frank and Ivanisvili used the analogous
limiting argument for the nearest-neighbor multiplier, which yields the log-Sobolev inequality for the 
heat-semigroup on \(\TT\) \cite{FrankIvanisvili2026}.
Indeed, after spectral-gap normalization, the nearest-neighbor multiplier
\[
  \ell_n(j)=\frac{1-\cos(2\pi j/n)}{1-\cos(2\pi/n)}
\]
satisfies \(\ell_n(k)\to k^2\) for each fixed \(k>0\), so its
limiting Dirichlet form is \(\sum_{k\in\ZZ}k^2\abs{\wh f(k)}^2\), corresponding
to the heat semigroup.  By contrast, the word-length multiplier
\(\psin(k)=\min\{k,n-k\}\) satisfies \(\psin(k)\to k\).   Thus the word-length limit yields the Poisson
Dirichlet form \(\sum_{k\in\ZZ}\abs{k}\,\abs{\wh f(k)}^2\).

Similar to the finite cyclic-group proof, the direct 3rd moment argument also
extends to the circle group case without any difficulty.  In
Section~\ref{sec:proofs} we provide the details. 

\section{The cubic reduction}



For any \(f:\Zn{n}\to \RR\), set
\[
  \En(f):=\Avg{f\Apsi f}
  =\sum_{k=0}^{n-1}\psin(k)\abs{\wh f(k)}^2.\]
We recall the cubic majorant used in
Frank and Ivanisvili's cubic Sobolev reduction \cite{FrankIvanisvili2026}.

\begin{lemma}[Cubic majorant {\cite{FrankIvanisvili2026}}]\label{lem:cubic-majorant}
For every \(t\ge0\),
\[
  2t^2\log t
  \le
  \frac23(t-1)^2(t+2)+(t^2-1),
\]
where the left side is interpreted as \(0\) at \(t=0\).
\end{lemma}

\begin{proposition}\label{prop:cubic-criterion}
Suppose that for every \(x:\Zn{n}\to[0,\infty)\) with \(\Avg{x^2}=1\),
\begin{equation}\label{eq:cubic-criterion}
  \En(x)\ge \frac13\Avg{(x-1)^2(x+2)}.
\end{equation}
Then
\[
  \Ent{f^2}\le 2\En(f)
\]
for every nonnegative \(f:\Zn{n}\to\RR\).
\end{proposition}

\begin{proof}
By homogeneity, normalize \(\Avg{f^2}=1\).  Then
\[
  \Ent{f^2}=2\Avg{f^2\log f}.
\]
Applying Lemma \ref{lem:cubic-majorant} pointwise and averaging gives
\[
  \Ent{f^2}
  \le \frac23\Avg{(f-1)^2(f+2)}
  +\Avg{f^2-1}
  =\frac23\Avg{(f-1)^2(f+2)}.
\]
The conclusion follows from \eqref{eq:cubic-criterion}.
\end{proof}

The next proposition isolates the estimate needed for the cubic criterion \eqref{eq:cubic-criterion}.

\begin{proposition}\label{prop:block-to-cubic}
Fix \(n\ge4\).  Assume that every real mean-zero \(u:\Zn{n}\to\RR\)
satisfies
\begin{equation}\label{eq:block-estimate-abstract}
  \abs{\Avg{u^3}}
  \le 3\Ssq\Qsymbol,
  \qquad
  \Ssq=\Avg{u^2},
  \quad
  \Qsq=\En(u)-\Ssq,
  \quad
  \Ssymbol,\Qsymbol\ge0.
\end{equation}
Then \eqref{eq:cubic-criterion} holds on \(\Zn{n}\).
\end{proposition}

\begin{proof}
Let \(x:\Zn{n}\to[0,\infty)\) and normalize \(\Avg{x^2}=1\).  Put
\[
  a=\Avg{x},\qquad u=x-a.
\]
Then \(\Avg{u}=0\), \(0\le a\le1\), and
\[
  \Ssq=\Avg{u^2}=1-a^2,
  \qquad
  \En(x)=\En(u)=\Ssq+\Qsq.
\]
Since \((t-1)^2(t+2)=t^3-3t+2\),
\[
  \Avg{(x-1)^2(x+2)}=\Avg{x^3}-3a+2.
\]
Expanding \(x=a+u\) gives
\[
  \Avg{x^3}=a^3+3a\Ssq+\Avg{u^3}
  =a^3+3a(1-a^2)+\Avg{u^3}.
\]
Hence
\begin{equation}\label{eq:cubic-expansion}
  \Avg{(x-1)^2(x+2)}=2-2a^3+\Avg{u^3}.
\end{equation}
Therefore \eqref{eq:cubic-criterion} is equivalent to
\begin{equation}\label{eq:needed-after-expansion}
  \Avg{u^3}\le (1-a)^2(1+2a)+3\Qsq.
\end{equation}
By \eqref{eq:block-estimate-abstract},
\[
  \Avg{u^3}\le 3\Ssq\Qsymbol.
\]
Also
\[
  3\Ssq\Qsymbol\le \frac34\Sfour+3\Qsq,
\]
and, since \(\Ssq=1-a^2\),
\[
  (1-a)^2(1+2a)-\frac34\Sfour
  =(1-a)^2(1+2a)-\frac34(1-a^2)^2
  =\frac14(1-a)^3(1+3a)\ge0.
\]
Thus \eqref{eq:needed-after-expansion} holds.
\end{proof}

\section{The 3rd moment estimate for odd cyclic groups}

Assume \(n=2m+1\ge5\).  For a real mean-zero function
\(u:\Zn{n}\to\RR\), write
\[
  u=\sum_{j=1}^{m}\left(\coef_j\ch{j}+\overline{\coef_j}\ch{-j}\right),
\]
where \(\ch{-j}=\ch{n-j}\).
Define
\[
  a_j=\sqrt2\,\abs{\coef_j},\qquad 1\le j\le m.
\]
Then
\[
  \Ssq=\Avg{u^2}=\sum_{j=1}^{m}a_j^2,
  \qquad
  \Qsq=\sum_{j=2}^{m}(j-1)a_j^2,
  \qquad \Ssymbol,\Qsymbol\ge0.
\]

\begin{proposition}[Odd 3rd moment estimate]\label{thm:odd-block}
For every real mean-zero \(u:\Zn{2m+1}\to\RR\), \(m\ge2\),
\[
  \abs{\Avg{u^3}}
  \le 3\Ssq\Qsymbol.
\]
\end{proposition}

\begin{proof}
If \(\Qsymbol=0\), then \(u\) is supported on the first frequency block.  Since
\(n\ge5\) is odd, no three elements of \(\{1,-1\}\) sum to zero modulo
\(n\).  Hence \(\Avg{u^3}=0\).  Assume \(\Qsymbol>0\).

The zero-frequency cubic resonances are precisely
\[
  i+j=k,
  \qquad
  i+j+k=n,
  \qquad 1\le i,j,k\le m,
\]
together with conjugate sign patterns.  For \(2\le k\le m\), put
\[
  L_k=\sum_{\substack{i+j=k\\1\le i,j\le m}}a_i a_j,
  \qquad L_1=0,
\]
and for \(1\le k\le m\), put
\[
  R_k=\sum_{\substack{i+j=n-k\\1\le i,j\le m}}a_i a_j.
\]
Coefficient counting gives
\begin{equation}\label{eq:odd-resonance-bound}
  \abs{\Avg{u^3}}
  \le \frac1{\sqrt2}\sum_{k=1}^{m}a_k(3L_k+R_k).
\end{equation}
Indeed, a product corresponding to \(i+j=k\) has magnitude
\(a_i a_j a_k/(2\sqrt2)\); the three placements of the negative
frequency and the conjugate sign pattern give the contribution
\(3a_kL_k/\sqrt2\).  For \(i+j+k=n\), the ordered all-positive triples
and the all-negative conjugates give \(a_kR_k/\sqrt2\).

The boundary term \(k=1\) is not covered by the \(\Qsq\)-weight.  Since
\(R_1\) corresponds to \(i+j=2m\), the only ordered pair is \((m,m)\), and
\(R_1=a_m^2\).  Thus
\[
  B_0=\frac1{\sqrt2}a_1R_1
  =\frac1{\sqrt2}a_1a_m^2.
\]
Using
\[
  \Ssq\ge a_1^2+a_m^2\ge2a_1a_m,
  \qquad
  \Qsq\ge(m-1)a_m^2,
\]
we get
\begin{equation}\label{eq:odd-boundary}
  B_0\le \frac1{2\sqrt{2(m-1)}}\Ssq\Qsymbol
  \le \frac1{2\sqrt2}\Ssq\Qsymbol.
\end{equation}

It remains to estimate
\[
  B_1=\frac1{\sqrt2}\sum_{k=2}^{m}a_k(3L_k+R_k).
\]
By Cauchy--Schwarz with weight \(k-1\),
\[
  B_1\le \frac1{\sqrt2}\Qsymbol\left(
  \sum_{k=2}^{m}\frac{(3L_k+R_k)^2}{k-1}
  \right)^{1/2}.
\]
We claim that
\begin{equation}\label{eq:odd-W}
  \sum_{k=2}^{m}\frac{(3L_k+R_k)^2}{k-1}
  \le 11\Sfour.
\end{equation}
For \(2\le k\le m\), set
\[
  U_k=\sum_{\substack{i+j=k\\1\le i,j\le m}}a_i^2 a_j^2,
  \qquad
  V_k=\sum_{\substack{i+j=n-k\\1\le i,j\le m}}a_i^2 a_j^2.
\]
The sum \(L_k\) has \(k-1\) ordered terms, while \(R_k\) has \(k\)
ordered terms.  Hence
\[
  L_k^2\le(k-1)U_k,
  \qquad
  R_k^2\le kV_k.
\]
Since \(k/(k-1)\le2\),
\[
  \frac{(3L_k+R_k)^2}{k-1}
  \le \left(3\sqrt{U_k}+\sqrt2\sqrt{V_k}\right)^2
  \le 11(U_k+V_k).
\]
The \(U_k\) terms cover ordered pairs with \(2\le i+j\le m\), and the
\(V_k\) terms cover ordered pairs with \(m+1\le i+j\le2m-1\).  These
ranges are disjoint and contained in all ordered pairs \(1\le i,j\le m\).
Thus
\[
  \sum_{k=2}^{m}(U_k+V_k)
  \le \sum_{i,j=1}^{m}a_i^2 a_j^2=\Sfour,
\]
which proves \eqref{eq:odd-W}.  Therefore
\begin{equation}\label{eq:odd-B1}
  B_1\le \sqrt{\frac{11}{2}}\Ssq\Qsymbol.
\end{equation}
Combining \eqref{eq:odd-resonance-bound}, \eqref{eq:odd-boundary}, and
\eqref{eq:odd-B1},
\[
  \abs{\Avg{u^3}}
  \le \left(\sqrt{\frac{11}{2}}+\frac1{2\sqrt2}\right)\Ssq\Qsymbol
  <3\Ssq\Qsymbol.
\]
\end{proof}

\section{The 3rd moment estimate for even cyclic groups}

Let \(n=2m\ge4\).  The proof is parallel to the odd case, with one
change: the middle frequency \(m=n/2\) is self-conjugate.  Write a real
mean-zero function as
\[
  u=\sum_{j=1}^{m-1}\left(\coef_j\ch{j}+\overline{\coef_j}\ch{-j}\right)
  +b\ch{m},
  \qquad b\in\RR,
\]
where \(\ch{m}(x)=(-1)^x\).  Again, the paired frequencies are only conjugate pairs coming from
real-valuedness.  The middle frequency is its own negative modulo \(2m\),
so its coefficient \(b\) is real.  Put
\[
  a_j=\sqrt2\,\abs{\coef_j}\quad (1\le j\le m-1),
  \qquad
  a_m=\abs{b}.
\]
Then
\[
  \Ssq=\sum_{j=1}^{m}a_j^2,
  \qquad
  \Qsq=\sum_{j=2}^{m}(j-1)a_j^2,
  \qquad \Ssymbol,\Qsymbol\ge0.
\]

\begin{proposition}[Even 3rd moment estimate]\label{thm:even-block}
For every real mean-zero \(u:\Zn{2m}\to\RR\), \(m\ge2\),
\[
  \abs{\Avg{u^3}}
  \le 3\Ssq\Qsymbol.
\]
\end{proposition}

\begin{proof}
The case \(\Qsymbol=0\) is again trivial, because then only the first frequency
block is present and no cubic zero-frequency resonance occurs for \(n\ge4\).
Assume \(\Qsymbol>0\).

For \(2\le k\le m-1\), define
\[
  L_k=\sum_{\substack{i+j=k\\1\le i,j\le m-1}}a_i a_j,
  \qquad
  R_k=\sum_{\substack{i+j=2m-k\\1\le i,j\le m-1}}a_i a_j,
\]
and define the middle-frequency sum
\[
  L_m=\sum_{\substack{i+j=m\\1\le i,j\le m-1}}a_i a_j.
\]
Empty sums are understood as zero.  The ordinary and wrap-around resonances give the
same coefficients as in the odd case, while the self-conjugate middle
frequency contributes with coefficient \(3\).  Hence
\begin{equation}\label{eq:even-resonance-bound}
  \abs{\Avg{u^3}}
  \le
  \frac1{\sqrt2}\sum_{k=2}^{m-1}a_k(3L_k+R_k)
  +3a_mL_m.
\end{equation}
By Cauchy--Schwarz, one has
\[
  \abs{\Avg{u^3}}
  \le \Qsymbol\sqrt{W},
\]
where
\[
  W=
  \sum_{k=2}^{m-1}\frac{(3L_k+R_k)^2}{2(k-1)}
  +\frac{9L_m^2}{m-1}.
\]
We prove
\begin{equation}\label{eq:even-W}
  W\le9\Sfour.
\end{equation}
For \(2\le k\le m-1\), let
\[
  U_k=\sum_{\substack{i+j=k\\1\le i,j\le m-1}}a_i^2 a_j^2,
  \qquad
  V_k=\sum_{\substack{i+j=2m-k\\1\le i,j\le m-1}}a_i^2 a_j^2.
\]
Both \(L_k\) and \(R_k\) have \(k-1\) ordered terms, so
\[
  L_k^2\le(k-1)U_k,
  \qquad
  R_k^2\le(k-1)V_k.
\]
Therefore
\[
  \frac{(3L_k+R_k)^2}{2(k-1)}
  \le \frac12\left(3\sqrt{U_k}+\sqrt{V_k}\right)^2
  \le5(U_k+V_k),
\]
using \((3a+b)^2\le10(a^2+b^2)\).  For the middle-frequency term, put
\[
  U_m=\sum_{\substack{i+j=m\\1\le i,j\le m-1}}a_i^2 a_j^2.
\]
Since \(L_m\) has \(m-1\) ordered terms,
\[
  \frac{9L_m^2}{m-1}\le9U_m.
\]
The \(U_k\) terms cover sums \(2\le i+j\le m-1\), the \(U_m\) term covers
sum \(m\), and the \(V_k\) terms cover sums \(m+1\le i+j\le2m-2\).  These
families partition the ordered pairs \(1\le i,j\le m-1\).  Hence
\[
  \sum_{k=2}^{m-1}(U_k+V_k)+U_m
  =\left(\sum_{j=1}^{m-1}a_j^2\right)^2.
\]
Consequently
\[
  W\le
  5\left(\sum_{j=1}^{m-1}a_j^2\right)^2+4U_m
  \le9\left(\sum_{j=1}^{m-1}a_j^2\right)^2
  \le9\Sfour.
\]
This proves \eqref{eq:even-W}, and therefore
\[
  \abs{\Avg{u^3}}
  \le \Qsymbol\sqrt{9\Sfour}=3\Ssq\Qsymbol.
\]
\end{proof}


\section{Proofs of the main results}\label{sec:proofs}

\begin{proof}[Proof of Theorem \ref{thm:main}]
If \(n\ge4\) is odd, then \(n=2m+1\) with \(m\ge2\), and Proposition
\ref{thm:odd-block} verifies the hypothesis of Proposition
\ref{prop:block-to-cubic}.  If \(n\ge4\) is even, then Proposition
\ref{thm:even-block} verifies the same hypothesis.  Thus the cubic
criterion holds for all \(n\ge4\), and Proposition \ref{prop:cubic-criterion}
gives
\[
  \Ent{f^2}\le2\Avg{f\Apsi f}.
\]
The sharpness of the constant is a standard perturbation argument, since the spectral gap is 1. 
\end{proof}
% It remains to justify sharpness of the constant.  Let \(u\) be any nonzero
% real first-eigenfunction, for instance
% \(u(x)=\cos(2\pi x/n)\), so that \(\Avg{u}=0\) and \(\Apsi u=u\).  For
% \(\abs{\varepsilon}\) small, \(f_\varepsilon=1+\varepsilon u\) is
% nonnegative.  As \(\varepsilon\to0\),
% \[
%   \Ent{f_\varepsilon^2}=2\varepsilon^2\Avg{u^2}+O(\varepsilon^3),
%   \qquad
%   \Avg{f_\varepsilon\Apsi f_\varepsilon}
%   =\varepsilon^2\Avg{u\Apsi u}
%   =\varepsilon^2\Avg{u^2}.
% \]
% Hence the ratio of the entropy to the energy tends to \(2\), so no smaller
% constant can hold, already for this fixed \(n\).



We finish with a continuous-circle proof for completeness. 

\begin{proof}[Proof of Theorem \ref{thm:circle}]
% The finite-cycle proof is the standard Riemann-sum limit used by Yao for the
% word-length multiplier.  For a nonnegative trigonometric polynomial
% \(f\) of degree at most \(d\), sample \(f\) on \(\Zn{n}\) by
% \(f_n(x)=f(x/n)\), with \(n>\max\{2d,4\}\).  On these frequencies
% \(\psin(k)=\abs{k}\), so the finite-cycle Dirichlet form agrees with the
% circle Dirichlet form, while the entropies converge by Riemann sums.  Applying
% Theorem \ref{thm:main} and letting \(n\to\infty\) gives the polynomial case;
% the usual density and lower-semicontinuity argument gives the form-domain
% statement.  This proves the theorem by the finite-cycle route.

For completeness, we also record the direct 3rd moment estimate.  As in
Proposition \ref{prop:block-to-cubic}, it is enough to prove the circle analogue
of \eqref{eq:block-estimate-abstract}.  Let \(u\) be a real mean-zero
trigonometric polynomial and write
\[
  u=\sum_{j\ge1}\left(\coef_j\chi_j+\overline{\coef_j}\chi_{-j}\right),
  \qquad a_j=\sqrt2\,\abs{\coef_j}.
\]
Only finitely many of the \(a_j\)'s are nonzero.  Then
\[
  \Ssq=\AvgT{u^2}=\sum_{j\ge1}a_j^2,
  \qquad
  \Qsq=\AvgT{u\AT u}-\Ssq=\sum_{j\ge2}(j-1)a_j^2,
  \qquad \Ssymbol,\Qsymbol\ge0.
\]
If \(\Qsymbol=0\), only the first frequency block is present and no cubic
zero-frequency resonance occurs, so \(\AvgT{u^3}=0\).  Assume \(\Qsymbol>0\).
The only zero-frequency cubic resonances are \(i+j=k\), together with their
conjugates.  For \(k\ge2\), put
\[
  L_k=\sum_{i+j=k}a_i a_j,
  \qquad
  U_k=\sum_{i+j=k}a_i^2a_j^2,
  \qquad i,j\ge1.
\]
Coefficient counting gives
\[
  \abs{\AvgT{u^3}}
  \le \frac3{\sqrt2}\sum_{k\ge2}a_kL_k.
\]
By Cauchy--Schwarz, we get
\[
  \abs{\AvgT{u^3}}
  \le \frac3{\sqrt2}\Qsymbol
  \left(\sum_{k\ge2}\frac{L_k^2}{k-1}\right)^{1/2}.
\]
Since \(L_k\) has \(k-1\) ordered terms,
\[
  L_k^2\le (k-1)U_k.
\]
Therefore
\[
  \sum_{k\ge2}\frac{L_k^2}{k-1}
  \le \sum_{k\ge2}U_k
  =\sum_{i,j\ge1}a_i^2a_j^2=\Sfour.
\]
Thus
\[
  \abs{\AvgT{u^3}}
  \le \frac3{\sqrt2}\Ssq\Qsymbol
  \le 3\Ssq\Qsymbol.
\]
The cubic reduction gives \(\EntT{f^2}\le2\AvgT{f\AT f}\) for nonnegative
trigonometric polynomials \(f\), and this concludes the proof.
%Sharpness follows from the perturbation \(f=1+\varepsilon\cos(2\pi x)\) and \(\varepsilon\to0\).
\end{proof}

\begin{thebibliography}{JPPP17}

\bibitem[And02]{Andersson2002}
M. E. Andersson,
\emph{Beitrag zur Theorie des Poissonschen Integrals \"uber endlichen Gruppen},
Monatsh. Math. \textbf{134} (2002), no. 3, 177--190.

\bibitem[Bec75]{Beckner1975}
W. Beckner,
\emph{Inequalities in Fourier analysis},
Ann. of Math. (2) \textbf{102} (1975), no. 1, 159--182.


\bibitem[BJJ83]{BJJ1983}
W. Beckner, S. Janson and D. Jerison,
\emph{Convolution inequalities on the circle},
in \emph{Conference on harmonic analysis in honor of Antoni Zygmund}, Vol. I,
II, Wadsworth Math. Ser., Wadsworth, Belmont, CA, 1983, 32--43.


\bibitem[Bon70]{Bonami1970}
A. Bonami,
\emph{\'{E}tude des coefficients de Fourier des fonctions de \(L^p(G)\)},
Ann. Inst. Fourier (Grenoble) \textbf{20} (1970), no. 2, 335--402.

\bibitem[CS03]{ChenSheu2003}
G.-Y. Chen and Y.-C. Sheu,
\emph{On the log-Sobolev constant for the simple random walk on the \(n\)-cycle:
the even cases},
J. Funct. Anal. \textbf{202} (2003), no. 2, 473--485.

\bibitem[CLS08]{ChenLiuSaloffCoste2008}
G.-Y. Chen, W.-W. Liu and L. Saloff-Coste,
\emph{The logarithmic Sobolev constant of some finite Markov chains},
Ann. Fac. Sci. Toulouse Math. (6) \textbf{17} (2008), no. 2, 239--290.

\bibitem[DS96]{DiaconisSaloffCoste1996}
P. Diaconis and L. Saloff-Coste,
\emph{Logarithmic Sobolev inequalities for finite Markov chains},
Ann. Appl. Probab. \textbf{6} (1996), no. 3, 695--750.

\bibitem[FF24]{FaustFawzi2024}
O. Faust and H. Fawzi,
\emph{Sum-of-squares proofs of logarithmic Sobolev inequalities on finite Markov chains},
IEEE Trans. Inform. Theory \textbf{70} (2024), no. 2, 803--819.


\bibitem[FI26]{FrankIvanisvili2026}
R. L. Frank and P. Ivanisvili,
\emph{The sharp log Sobolev inequality on finite cycles},
arXiv:2605.29035, 2026.

\bibitem[Gro75]{Gross1975}
L. Gross,
\emph{Logarithmic Sobolev inequalities},
Amer. J. Math. \textbf{97} (1975), no. 4, 1061--1083.

\bibitem[Gro06]{Gross2006}
L. Gross,
\emph{Hypercontractivity, logarithmic Sobolev inequalities, and applications: a
survey of surveys},
in \emph{Diffusion, quantum theory, and radically elementary mathematics},
Math. Notes \textbf{47}, Princeton Univ. Press, Princeton, NJ, 2006, 45--73.



\bibitem[JPPP17]{JPPP2017}
M. Junge, C. Palazuelos, J. Parcet and M. Perrin,
\emph{Hypercontractivity in group von Neumann algebras},
Mem. Amer. Math. Soc. \textbf{249} (2017), no. 1183, xii+83.







\bibitem[Wei80]{Weissler1980}
F. B. Weissler,
\emph{Logarithmic Sobolev inequalities and hypercontractive estimates on the circle},
J. Funct. Anal. \textbf{37} (1980), no. 2, 218--234.

\bibitem[Wol07]{Wolff2007}
P. Wolff,
\emph{Hypercontractivity of simple random variables},
Studia Math. \textbf{180} (2007), no. 3, 219--236.

\bibitem[Yao25]{Yao2025}
G. Yao,
\emph{Optimal hypercontractivity and log-Sobolev inequalities on cyclic groups
\(\ZZ_{m\cdot2^k}\)},
arXiv:2512.03489, 2025.

\bibitem[Yao26]{Yao2026}
G. Yao,
\emph{A complete solution of optimal hypercontractivity and Log--Sobolev
inequalities on cyclic groups \(\ZZ_n\) for \(n\ge4\)},
arXiv:2605.30867v1, 2026.

\end{thebibliography}


\end{document}
